% =================================================
% Layout LaTeX per Header e Footer
% =================================================
%
% Questo file definisce l'header e il footer.
% Viene popolato dallo script generate-pdf.ps1
% e poi incluso nel file .tex generato da Pandoc.

% --- Configurazione Margini ---
% Questi valori sono già impostati in Pandoc, ma li ridefiniamo per sicurezza
\geometry{a4paper, top=3.5cm, bottom=2.5cm, left=2cm, right=2cm, headheight=3cm, footskip=1.5cm}

% --- Definizione Stile Pagina "fancy" ---
\pagestyle{fancy}
\fancyhf{} % Pulisce tutti i campi di header e footer

\setlength{\arrayrulewidth}{1pt}

% --- Definizione colori personalizzati ---
\definecolor{azzurro}{HTML}{5F96DD}
\definecolor{arancio}{HTML}{9D5700}
\definecolor{viola}{HTML}{898DC3}

% --- Configurazione Header ---
\fancyhead[L]{
    \begin{tikzpicture}
        \node[opacity=0.6, anchor=north west] at (0,0) {
            \includegraphics[height=1.5cm]{
                % --- Comment line to prevent powershell substitution issue
                {{LOGO_PATH}}
            }
        };
    \end{tikzpicture}%
}

\newcolumntype{C}[1]{>{\centering\arraybackslash}p{#1}}

\fancyhead[R]{%
    \begin{tikzpicture}[remember picture, overlay]

        % Nodo contenente la tabella
        \node[opacity=0.6, anchor=south east, yshift=2.0cm, inner sep=0pt] (tabella_header) at (0.3,-1.8) {
            {\small
                \color{arancio}
                \renewcommand{\arraystretch}{0.8}
                \begin{tabular}{ |C{7cm}|C{4cm}|C{1.4cm}| }
                    \hline
                    \color{black}\textit{Politica}\color{arancio} & 
                    \color{black}\textit{Codice Documento}\color{arancio} & 
                    \color{black}\textit{Versione}\color{arancio} \\ 
                    \color{black}{{HEADER_POLITICA}}\color{arancio} &
                    \color{black}{{HEADER_CODICE}}\color{arancio} &
                    \color{black}{{HEADER_VERSIONE}}\color{arancio}
                \end{tabular}
            }
        };

        % Linea viola sotto la tabella usando coordinate del nodo
        \draw[viola, line width=4pt, opacity=0.6] 
            ([xshift=4pt,yshift=-1pt]tabella_header.south west) -- 
            ([xshift=-3pt,yshift=-1pt]tabella_header.south east);

    \end{tikzpicture}%
}

% --- Configurazione Footer ---
\fancyfoot[L]{
    \begin{tikzpicture}[remember picture, overlay]
        \node[opacity=0.6, anchor=south west, yshift=-0.4cm] at (0,0) {
            % --- Comment line to prevent powershell substitution issue
            {{FILENAME}}
        };
    \end{tikzpicture}
}
        
\fancyfoot[R]{
    \begin{tikzpicture}[remember picture, overlay]
        \node[opacity=0.6, anchor=south east, yshift=-0.4cm] at (0.3,0) {
            \thepage
        };
    \end{tikzpicture}
}

% --- Linea di separazione ---
\renewcommand{\headrulewidth}{0pt}
\renewcommand{\footrulewidth}{0pt}

% --- Applica lo stile di pagina ---
% Questo comando imposta lo stile "fancy" per tutte le pagine
% La copertina userà \thispagestyle{empty} per non avere header/footer.
\pagestyle{fancy}